\begin{afterword}
\summaryhead

The author recounts the moment, in late 2002, of seeing what the AI proposals of 1997 would really have done. Two reactions are refused: the proud ``I almost destroyed the world!'', and the claim that the gun was not really loaded, since the author had intended to build it.

The lesson was to stop and say ``I'm not ready,'' words that the author says cost a great deal of status in AGI research. The deeper reason the author did not stop sooner was not believing at gut level that ``Nature was allowed to kill me.'' The post guesses that other AGI researchers, who fear being beaten by rivals, are in the same state.

The author then understood that there is no grading on a curve, and that good arguments for a risk do not stop it from killing you. The only safe course in ignorance was not to proceed until the ground was known to be safe. That, the post says, is when the author's training as a rationalist began.

\responsehead

The post is memoir, and the first question is whether its account of the author's earlier views and writings holds up. Where it can be checked, it does. In 1996 the author wrote on the Extropians mailing list that nanotechnology would turn the Earth into gray goo, and in the same message listed whether the post-Singularity ``Powers will be ethical'' among questions answered by ``Almost certainly'', ``of course'', and ``consider the alternative.'' In 1999 the author planned an open-source AI effort ``which I intend to oversee at some point.'' \textsc{Creating Friendly AI} (2001) argued that delay would not make Friendliness easier, and that many harmless mistakes would show up before a catastrophic one. The post calls these ``careful arguments'' for going on at full speed and reasons ``to tolerate the risk,'' and the texts fit the description.

The post is also careful with its own claims. It limits what the 1997 proposals ``really would have done'' to what was coherent enough to judge. It refuses both the proud and the minimizing account of the mistake. It marks its reading of other researchers as a guess: ``It is unsafe to say what other people are thinking.''

The weaker parts are the claims about the field. That saying ``I'm not ready to write code'' costs great status among AGI researchers and the public is drawn from the author's experience; no case is given, and the reader is instead challenged to name a counterexample. The guess that the prospect of destroying humanity by their own mistake is not ``real'' to other researchers rests on the author's own past: ``I, too, used to say things like that.''

The post quotes the argument that a project which delays for safety will be beaten by projects that do not care. Its answer is to set rivals aside: even the best project can be killed. It does not say what stopping achieves if the others go on, which is the case that argument raises. In 2001 the author had used that case as a reason for speed.

In short: a memoir of the author's 2002 turn that the archived writings bear out, with claims about the AGI field drawn from the author's experience, and an answer to the race argument that sets rivals aside without saying what stopping achieves while they go on.
\end{afterword}
